% Дүгнэлт -- 4-р бүлэгийн үргэлжлэл (Судалгааны үр дүн, дүгнэлт)
\section{Дүгнэлт}
\label{sec:conclusion}

Энэхүү судалгааны ажлын хүрээнд ABU Робокон 2026 тэмцээний R2 бүрэн автомат роботод зориулсан объект таних систем болон шийдвэр гаргах алгоритмыг хөгжүүлж, симуляцийн орчинд туршин баталгаажуулсан болно. Судалгааны үндсэн үр дүнг дараах байдлаар нэгтгэн дүгнэж болно.

\subsection{Симуляцийн орчин}

Unity видео тоглоомын хөдөлгүүрт тулгуурлан тэмцээний талбайн бодитой симуляцийн орчныг бүрдүүлсэн. FreeCAD болон Blender программуудад боловсруулсан 3D загваруудыг FBX форматаар импортлон, физикийн харилцан үйлчлэлийг \textit{Collider} бүрэлдэхүүнээр тодорхойлсон. Энэхүү орчин нь зөвхөн туршилтын зориулалттай бус, датасет боловсруулахад мөн амжилттай ашиглагдсан.

\subsection{Датасет боловсруулалт}

Симуляцийн орчинд камерыг санамсаргүй байрлал, өнцөг, зайд байршуулах алгоритмаар 10,000 зургаас бүрдсэн датасетийг автоматаар үүсгэсэн. Зураг бүрт харгалзах тайлбар файлыг объектуудын дэлхийн координатыг камерын координатын системд хөрвүүлэх замаар тооцоолж үүсгэсэн. Датасет нь олон янзын гэрэлтүүлэг, өнцөг, байрлалын нөхцлийг хамарсан бөгөөд, нийт 30 ангиллын объектыг агуулсан.

\subsection{Объект таних алгоритм}

YOLOv12 загварыг шилжүүлэн сургах \textit{(transfer learning)} аргаар боловсруулсан датасет дээр сургасан. Сургалтын үр дүнд recall(B) хэмжүүр 91\%-д хүрсэн бөгөөд, симуляцийн орчинд хийсэн туршилтаар ихэнх объектыг 90\%-иас дээш итгэлтэйгээр илрүүлсэн байна. Буруу илрүүлэлт \textit{(false positive)} туршилтын явцад огт ажиглагдаагүй нь загварын найдвартай байдлыг харуулж байна.

Сургасан загварыг симуляцийн орчинтой холбохдоо Python хэл дээр HTTP сервер үүсгэж, түүхий дата \textit{(raw binary data)} форматаар дамжуулснаар JSON форматтай харьцуулахад 10 дахин бага хэмжээний дата солилцох боломжтой болсон.

\subsection{Шийдвэр гаргах алгоритм}

Роботын үйл ажиллагааг хязгаарлагдмал төлөвт автомат \textit{(Finite State Machine)} аргаар загварчилсан. Тохиргоо, зэвсэг угсрах, судар авах, судар байршуулах гэсэн 4 үндсэн төлөвтэй диаграм боловсруулсан бөгөөд, тэмцээний даалгаврын дарааллыг тодорхой, удирдахад хялбар байдлаар илэрхийлсэн.

\subsection{Цаашдын судалгааны чиглэл}

Энэхүү судалгааны ажлыг цаашид дараах чиглэлүүдээр хөгжүүлэх боломжтой:

\begin{itemize}
    \item \textbf{Бодит орчинд туршилт} --- Симуляцийн орчин болон бодит орчны ялгааг \textit{(sim-to-real gap)} судалж, загварын бодит нөхцөлд ажиллах чадварыг сайжруулах.
    \item \textbf{Датасетийн өргөтгөл} --- Датасетийн хэмжээ болон олон янз байдлыг нэмэгдүүлэх, ялангуяа далдлагдсан болон хэсэгчлэн харагдах объектуудын нөхцлийг илүүтэй хамруулах.
    \item \textbf{Гүйцэтгэлийн оновчлол} --- Объект таних алгоритмын боловсруулалтын хурдыг бодит цагийн шаардлагад нийцүүлэн оновчлох.
\end{itemize}
